% Mathématiques 5e — cours V3
% Dépendance entre grandeurs : première approche des fonctions

\section{Objectifs}

\begin{itemize}
\item Employer l'expression « en fonction de » dans des contextes concrets ou mathématiques.
\item Produire un tableau de valeurs.
\item Lire et interpréter un tableau de valeurs.
\item Placer dans un repère des points issus d'un tableau de valeurs.
\item Lire et interpréter un graphique cartésien donné par une courbe ou un nuage de points.
\item Produire un tableau de valeurs à partir d'une formule.
\item Produire une formule simple représentant la dépendance entre deux grandeurs.
\item Caractériser graphiquement la proportionnalité.
\end{itemize}

\section{1. Une grandeur dépend d'une autre}

De nombreuses situations relient deux grandeurs.

Par exemple :
\begin{itemize}
\item le prix dépend du nombre d'objets achetés ;
\item l'aire d'un carré dépend de la longueur de son côté ;
\item la distance parcourue dépend du temps ;
\item la température peut dépendre de l'heure de la journée.
\end{itemize}

On dit alors qu'une grandeur est exprimée \textbf{en fonction de} l'autre.

\coursefigure{/images/mathematiques/5eme/fonctions-dependance.svg}{Une grandeur d'entrée passe par une formule ou une relation pour produire une grandeur de sortie.}{Dire qu'une grandeur dépend d'une autre signifie que sa valeur est déterminée à partir de la première selon une règle donnée.}

\section{2. Exemple concret}

Un parking coûte :
\[
2\ \text{€}
\]
par heure.

Le prix dépend de la durée.

Pour une durée $t$ en heures :
\[
P=2t.
\]

On peut dire :
\[
\text{« le prix }P\text{ est exprimé en fonction de }t\text{ »}.
\]

\section{3. Ce qui change et ce qui dépend}

Dans :
\[
A=c^2,
\]
l'aire $A$ d'un carré dépend de la longueur $c$ de son côté.

Si :
\[
c=3\ \text{cm},
\]
alors :
\[
A=3^2=9\ \text{cm}^2.
\]

Si :
\[
c=5\ \text{cm},
\]
alors :
\[
A=5^2=25\ \text{cm}^2.
\]

Une même règle produit différentes valeurs selon l'entrée.

\section{4. Tableau de valeurs}

Un tableau de valeurs associe plusieurs valeurs de la grandeur de départ aux valeurs correspondantes.

Pour :
\[
P=2t,
\]
on peut construire :

\begin{tabular}{c|c|c|c|c}
$t$ en h & $0$ & $1$ & $2$ & $3$ \\
$P$ en € & $0$ & $2$ & $4$ & $6$ \\
\end{tabular}

Chaque colonne représente un couple de valeurs liées par la même règle.

\section{5. Produire un tableau à partir d'une formule}

Considérons :
\[
y=3x+1.
\]

Pour :
\[
x=0,
\]
on obtient :
\[
y=3\times0+1=1.
\]

Pour :
\[
x=2,
\]
on obtient :
\[
y=3\times2+1=7.
\]

On peut construire :

\begin{tabular}{c|c|c|c|c}
$x$ & $0$ & $1$ & $2$ & $3$ \\
$y$ & $1$ & $4$ & $7$ & $10$ \\
\end{tabular}

\section{6. Lire un tableau}

Dans le tableau précédent, à :
\[
x=2
\]
correspond :
\[
y=7.
\]

On peut aussi poser la question inverse :
\[
\text{« pour quelle valeur de }x\text{ obtient-on }y=10\text{ ? »}
\]

La lecture du tableau donne :
\[
x=3.
\]

\section{7. Tableau incomplet}

Un tableau peut comporter des cases à compléter.

Si la règle est :
\[
y=4x,
\]
et que :
\[
x=2{,}5,
\]
alors :
\[
y=4\times2{,}5=10.
\]

Si :
\[
y=28,
\]
alors :
\[
4x=28,
\]
donc :
\[
x=7.
\]

Le tableau peut ainsi être lu dans les deux sens.

\section{8. Du tableau au repère}

Chaque colonne d'un tableau peut devenir un point du plan.

Par exemple :
\[
x=1,\quad y=4
\]
donne le point :
\[
(1;4).
\]

\coursefigure{/images/mathematiques/5eme/fonctions-tableau.svg}{Tableau de valeurs relié à plusieurs points placés dans un repère.}{Chaque couple du tableau devient un point ; le graphique permet de visualiser la dépendance entre les grandeurs.}

\section{9. Placer les points}

Pour le tableau :
\[
(0;1),\quad(1;4),\quad(2;7),\quad(3;10),
\]
on place quatre points dans un repère.

Il faut :
\begin{enumerate}
\item lire l'abscisse ;
\item lire l'ordonnée ;
\item respecter les graduations des axes ;
\item conserver les unités si les axes représentent des grandeurs physiques.
\end{enumerate}

\section{10. Lire un graphique}

Un graphique cartésien peut être constitué :
\begin{itemize}
\item de points isolés ;
\item d'un nuage de points ;
\item d'une courbe.
\end{itemize}

Pour lire la valeur correspondant à une abscisse :
\begin{enumerate}
\item partir de l'abscisse ;
\item se déplacer verticalement jusqu'au graphique ;
\item lire ensuite l'ordonnée correspondante.
\end{enumerate}

\section{11. Interpréter un graphique}

Un graphique ne représente pas seulement des nombres.

Il peut indiquer :
\begin{itemize}
\item une augmentation ;
\item une diminution ;
\item une valeur maximale ;
\item une valeur minimale ;
\item une stabilisation ;
\item des points particuliers.
\end{itemize}

L'interprétation doit toujours revenir au sens des axes.

\section{12. Exemple de température}

On relève la température toutes les deux heures.

Un graphique peut montrer que la température :
\begin{itemize}
\item augmente le matin ;
\item atteint un maximum dans l'après-midi ;
\item diminue ensuite.
\end{itemize}

Une lecture graphique reste approximative si les points exacts ne sont pas indiqués.

\section{13. Produire une formule}

Une formule doit traduire la règle de dépendance.

\begin{exemple}
Un service facture :
\[
5\ \text{€}
\]
de frais fixes et :
\[
3\ \text{€}
\]
par unité consommée.

Si $x$ est la quantité consommée, le prix $P$ vaut :
\[
\boxed{P=3x+5}.
\]
\end{exemple}

La formule résume tous les cas possibles.

\section{14. Formule géométrique}

Le périmètre d'un rectangle de longueur :
\[
L
\]
et de largeur fixe :
\[
3
\]
vaut :
\[
P=2L+2\times3.
\]

Donc :
\[
\boxed{P=2L+6}.
\]

Le périmètre est ici exprimé en fonction de $L$.

\section{15. Passer de la formule au tableau}

Avec :
\[
P=2L+6,
\]
on choisit plusieurs valeurs de $L$.

\begin{tabular}{c|c|c|c}
$L$ & $1$ & $3$ & $5$ \\
$P$ & $8$ & $12$ & $16$ \\
\end{tabular}

Le tableau permet ensuite de placer les points correspondants.

\section{16. Dépendance proportionnelle}

Une proportionnalité est un cas particulier de dépendance.

Si :
\[
y=4x,
\]
les valeurs sont proportionnelles.

Le graphique est constitué de points alignés avec l'origine.

\begin{propriete}
Une relation proportionnelle se reconnaît graphiquement par des points alignés avec :
\[
O(0;0).
\]
\end{propriete}

\section{17. Dépendance non proportionnelle}

Avec :
\[
y=4x+3,
\]
on a :
\[
x=0\Rightarrow y=3.
\]

Le graphique ne passe donc pas par l'origine.

La relation exprime une dépendance, mais elle n'est pas proportionnelle.

\section{18. Exemple comparatif}

Situation A :
\[
P=2x.
\]

Situation B :
\[
Q=2x+5.
\]

Dans la situation A :
\[
x=0\Rightarrow P=0.
\]

Dans la situation B :
\[
x=0\Rightarrow Q=5.
\]

La première peut être proportionnelle.

La seconde comporte une valeur fixe et ne l'est pas.

\section{19. Données discrètes}

Certaines situations ne donnent que quelques valeurs.

Par exemple, le nombre de billets achetés est entier :
\[
0,\quad1,\quad2,\quad3,\ldots
\]

Le graphique peut alors être un ensemble de points non reliés.

Relier automatiquement les points pourrait suggérer des valeurs qui n'ont pas de sens dans le contexte.

\section{20. Données continues}

Une grandeur comme le temps peut prendre des valeurs intermédiaires.

Dans certains contextes, une courbe continue peut alors être adaptée.

Le type de graphique dépend donc aussi de la nature des données.

\section{21. Exemple développé}

La distance parcourue à vitesse constante :
\[
60\ \text{km/h}
\]
est donnée par :
\[
d=60t.
\]

Pour :
\[
t=0{,}5\ \text{h},
\]
on obtient :
\[
d=60\times0{,}5=30\ \text{km}.
\]

Pour :
\[
t=2\ \text{h},
\]
on obtient :
\[
d=120\ \text{km}.
\]

Le tableau est :

\begin{tabular}{c|c|c|c}
$t$ en h & $0$ & $0{,}5$ & $2$ \\
$d$ en km & $0$ & $30$ & $120$ \\
\end{tabular}

Les points sont alignés avec l'origine : il s'agit d'une proportionnalité.

\section{22. Méthode : formule vers tableau}

\begin{methode}
\begin{enumerate}
\item choisir les valeurs d'entrée ;
\item remplacer la lettre dans la formule ;
\item calculer la valeur correspondante ;
\item reporter le couple dans le tableau ;
\item contrôler les unités.
\end{enumerate}
\end{methode}

\section{23. Méthode : tableau vers graphique}

\begin{methode}
\begin{enumerate}
\item choisir les axes ;
\item indiquer les grandeurs et unités ;
\item choisir une graduation adaptée ;
\item placer chaque couple comme un point ;
\item relier les points uniquement si le contexte le justifie.
\end{enumerate}
\end{methode}

\section{24. Méthode de lecture graphique}

\begin{methode}
Pour lire une valeur :
\begin{enumerate}
\item identifier l'axe de départ ;
\item repérer la valeur donnée ;
\item rejoindre le graphique ;
\item projeter sur l'autre axe ;
\item lire la valeur avec l'unité.
\end{enumerate}
\end{methode}

\section{25. Limite du programme de 5e}

\begin{attention}
En 5e, on introduit la dépendance entre grandeurs et l'expression « en fonction de ».

On ne construit pas encore une théorie générale des fonctions avec tout le vocabulaire d'image, d'antécédent et de fonction affine qui sera approfondi plus tard.
\end{attention}

L'objectif est d'apprendre à passer entre :
\[
\boxed{\text{contexte}\leftrightarrow\text{formule}\leftrightarrow\text{tableau}\leftrightarrow\text{graphique}}.
\]

\section{26. Erreurs fréquentes}

\begin{itemize}
\item Inverser abscisse et ordonnée.
\item Produire un tableau sans utiliser la même règle pour toutes les colonnes.
\item Lire un graphique sans regarder les unités.
\item Relier des points alors que les valeurs intermédiaires n'ont pas de sens.
\item Confondre toute dépendance avec une proportionnalité.
\item Croire que des points simplement alignés suffisent pour la proportionnalité sans vérifier l'origine.
\item Inventer une formule à partir de trop peu d'informations.
\end{itemize}

\section{À retenir}

\begin{aretenir}
Dire qu'une grandeur est « en fonction de » une autre signifie qu'elle dépend de celle-ci selon une règle.

Cette relation peut être représentée par :
\[
\boxed{\text{une formule, un tableau de valeurs ou un graphique}}.
\]

Chaque couple :
\[
(x;y)
\]
du tableau peut être placé dans un repère.

Une situation de proportionnalité constitue un cas particulier :
\[
\boxed{\text{les points sont alignés avec l'origine}}.
\]

En 5e, l'essentiel est de savoir passer d'une représentation à une autre et interpréter la dépendance dans son contexte.
\end{aretenir}
